%Theory
\section{Literature Review}
The Vehicle Routing Problem (VRP) belongs to the class of NP-hard problems and represents one of the most extensively studied topics in logistics and transportation. Its basic version, first proposed by Dantzig and Ramser in 1959, assumes the optimization of vehicle routes serving a set of customers while minimizing the total travel cost \cite{Dantzig1959}. Since then, many extensions have been developed, incorporating constraints such as vehicle capacity (CVRP), service times, different vehicle types (HVRP), and time windows within which customers can be served — the so-called Vehicle Routing Problem with Time Windows (VRPTW) \cite{Toth2014, Golden2008}. The capacitated version with time windows (CVRPTW) better reflects the conditions of real-world transportation systems, where resource availability and time requirements play a crucial role.

Modern logistics systems operate under increasingly dynamic conditions — not only in terms of variable transportation times and fluctuating fleet sizes but also regarding demand and changing time window boundaries. For example, in grocery retail networks, fresh products must be delivered regularly and in appropriate quantities to avoid losses related to product shelf life. In the construction industry, material deliveries must be synchronized with work schedules, as any delay may lead to downtime and costly disruptions. These situations demonstrate that routing models must account for both variability and uncertainty in operational data.

Traditional VRPTW and CVRPTW models assume deterministic data — fixed demand, travel times, and rigid time windows \cite{surveyBraekers, RCVRPTWreviewYernar}. In practice, however, these parameters frequently change, limiting the effectiveness of methods based solely on deterministic data. Consequently, there is growing interest in stochastic models and approaches aimed at ensuring greater solution robustness against data variability.

On one hand, the literature on stochastic VRP variants has been developing intensively — reviews indicate numerous studies addressing random demand and stochastic travel times \cite{Oyola2018, Bomboi2022}. One example is the work by Pavone et al. \cite{Pavone2009}, where a VRPTW model accounts for random arrival of requests, with time window starts depending on when demand appears. However, in most such studies, time windows remain deterministic, even though in real-world systems they may shift dynamically.

On the other hand, in the field of robust optimization, research has been conducted to safeguard the planning process against data uncertainty. For example, Munari et al. \cite{Munari2018} proposed an RVRPTW model with uncertain demand and travel times, along with a branch-price-and-cut algorithm. Zhang and Shen \cite{Zhang2021} developed a distributionally robust model for the VRPTW, minimizing the risk of delays under unknown travel time distributions. Similarly, Shen et al. \cite{Shen2022} analyzed electric vehicle routing problems with uncertain demand and energy consumption dependent on load and time.

Despite the rich literature, significant research gaps remain. Many works treat random demand or uncertain travel times as stochastic variables, while few consider variability in time windows themselves. Bomboi et al. \cite{Bomboi2022} showed that in VRPTW with random travel times, uncertainty strongly affects route feasibility. Errico et al. \cite{Errico2016} analyzed stochastic service times but under deterministic time windows. An interesting extension was proposed by Spliet and Gabor \cite{Spliet2014}, introducing the TWAVRP problem where time window boundaries are decisions made before demand realization. However, the literature still lacks works combining variable, random time windows with stochastic demand in the capacitated CVRPTW variant. Scenario-based approaches, where numerous random instances are generated and solved separately to assess stability and method behavior under uncertainty, are also rarely encountered. One of the newer directions includes works utilizing contextual data — for example, Serrano et al. \cite{Serano2024} analyzed travel times dependent on environmental features.

The systematic review by Liu et al. \cite{Liu2023} indicates that approximately 86\% of studies on VRPTW employ approximate methods, mainly heuristics and metaheuristics, with the majority being tested only on deterministic data. This means that despite the widespread use of advanced search methods, research on VRPTW under uncertainty still represents a smaller portion of the literature.

In summary, despite notable progress, clear research gaps can be identified in the following areas:
\begin{itemize}
    \item Modeling random time window boundaries as variables generated from probability distributions.
    \item Combining stochastic time windows with random demand within a single model.
    \item Applying scenario-based approaches enabling analysis of stability, risk, and operational robustness.
\end{itemize}

This work presents a method aligned with the research direction on operational robustness and stochastic analysis. The approach does not constitute formal robust optimization but allows evaluation of algorithm behavior across a broad set of diverse scenarios, providing valuable insights on solution stability and quality. This work thus provides a foundation for further research, both comparative studies and the development of new methods for handling uncertainty in routing problems.
